% ===========================================================================
% TASTI 2026 poster -- Paper No. T0043 (poster, 0201 Earth Observation and
% Remote Sensing)
%
%   Size   : A0 portrait, 84.1 cm x 118.9 cm (official requirement, see
%            docs/poster_requirements.md)
%   Build  : `make`          -> poster.pdf (converts figures first)
%            `make preview`  -> preview/poster.png for a quick look
%            `make check`    -> page size + embedded fonts
%
% Every number on this sheet is traced to a benchmark series in
% docs/results.md. The three series (S1 solver duel on one GPU, S2 large
% scene vs CPU, S3 five-scene end-to-end) must never be mixed, and no size
% dependence may be claimed from S3. The S1 ratio (16.5x) is allowed here
% ONLY with its qualifiers (same GPU, Fernandina, MintPy-reported step time
% without interpreter start-up); its step wall (4.49x) is never quoted.
%
% Block-level content plan and the reasoning behind it: docs/poster_plan.md.
% ===========================================================================
\documentclass[final,t]{beamer}
\usepackage[orientation=portrait,size=a0,scale=1.15]{beamerposter}
\usetheme{tasti2026}

\usepackage[utf8]{inputenc}
\usepackage[T1]{fontenc}
\usepackage{lmodern}
\usepackage{amsmath,amssymb}
\usepackage{booktabs}
\usepackage{graphicx}
\usepackage{qrcode}
\usepackage{textcomp}
\usepackage{microtype}

\usefonttheme{professionalfonts}
\renewcommand{\familydefault}{\sfdefault}
\setlength{\parskip}{0.5ex}

\graphicspath{{figures/}}

% Column geometry: three equal columns with a small gutter.
\newlength{\colw}
\setlength{\colw}{0.305\paperwidth}
\newlength{\gutter}
\setlength{\gutter}{0.0175\paperwidth}

% Repository and upstream links (also encoded as QR codes).
\newcommand{\prurl}{https://github.com/insarlab/MintPy/pull/1490}
\newcommand{\posterurl}{https://github.com/s-sasaki-earthsea-wizard/tasti2026-poster}

% ---------------------------------------------------------------------------
% Title block (title, paper number and author names must be at the top)
% ---------------------------------------------------------------------------
\title{Accelerating Open-Source InSAR Time-Series Analysis on Consumer GPUs:
  A PyTorch-Based Network Inversion Solver for MintPy}
\author{Syota Sasaki\textsuperscript{*}}
\institute{Earthsea Wizard, Japan \quad\textbullet\quad
  \texttt{syota.sasaki@earthsea-wizard.com} \quad\textbullet\quad
  \textsuperscript{*}presenter}

\renewcommand{\tastiHeadlineRight}{%
  \begin{beamercolorbox}[wd=0.18\paperwidth,sep=0.8ex,center]{paperno in headline}%
    \usebeamerfont{paperno in headline}Paper No.\ T0043\\[0.2ex]
    {\normalfont\small Poster \textbullet\ 0201 Earth Observation and Remote Sensing}
  \end{beamercolorbox}}

\renewcommand{\tastiFootline}{%
  TASTI 2026, Tainan, Taiwan \textbullet\
  Poster session: 10 November 2026, 11:30--14:00, Room Tayouan D (3F)
  \hfill
  Sources and benchmark reports: \texttt{github.com/s-sasaki-earthsea-wizard/tasti2026-poster}}

\begin{document}
\begin{frame}[t]
\begin{columns}[t]

% ===========================================================================
% COLUMN 1 -- problem
% ===========================================================================
\begin{column}{\colw}

% Block 1: the hotspot and the headline result. 43.5 % is report_baseline.md;
% 44x is S2 internal (same number as the block 6 big number). Earthquakes
% appear only as slow inter/post-seismic motion: co-seismic offsets are a
% single-interferogram problem, not a time-series one.
\begin{block}{1 \ The bottleneck}
  InSAR time-series analysis is a standard tool for tracking slow ground
  deformation over months to years: volcanic inflation, motion along active
  faults before and after earthquakes, and land subsidence. With
  Sentinel-1 acquiring continuously and NISAR~[1] now in orbit, archives
  grow faster than the throughput of the software that consumes them.

  \bignumber{43.5\,\%}{of a complete MintPy~[2] run (18-step
    \texttt{smallbaselineApp}) is spent in \emph{one} step: the least-squares
    inversion of the interferogram network, solved pixel by pixel on the CPU}

  \keypoint{This work: that step now runs 44$\times$ faster on one consumer
    GeForce GPU (block~6), switched on by a single template option.}
\end{block}

% Block 2: why the inversion exists at all, for readers new to InSAR. No
% benchmark numbers. Redundancy averages per-interferogram noise only:
% atmosphere belongs to a date, so it survives the inversion. The last bullet
% (per-pixel weights forbid a shared factorization) is the bridge to block 3.
\begin{block}{2 \ Why a least-squares inversion?}
  \begin{itemize}
    \item A radar satellite images the same ground on a fixed repeat cycle
      (12~days for Sentinel-1 and NISAR). An \emph{interferogram}, the phase
      difference of two images, measures how far the ground moved toward the
      satellite between those two dates, with millimetre-scale sensitivity.
    \item Pairs far apart in time or in orbit lose coherence, so the time
      series is built from many short pairs: a \emph{network} with dates as
      vertices, interferograms as edges, and more edges than unknown dates.
  \end{itemize}

  \centering
  \includegraphics[width=0.6\linewidth]{observation_network_example.pdf}
  \figcaption{A toy network: four dates and five interferograms. Each edge
    observes the displacement between its two dates; real networks have tens
    to hundreds of dates.}

  \raggedright
  \begin{itemize}
    \item Recovering the displacement at every date is then an overdetermined
      linear problem, solved by least squares; the redundancy averages out
      the noise of individual interferograms.
    \item Interferogram quality differs from pixel to pixel, so each pixel
      weights its equations differently. With shared weights, one
      factorization would serve every pixel; per-pixel weights turn the
      inversion into millions of separate small problems.
  \end{itemize}
\end{block}

% Block 3: the shape of the problem. Pixel counts: S2 (3.4 M) and
% compute_scale.md (NISAR frame). D^3 growth: compute_scale.md.
\begin{block}{3 \ The problem: millions of small least-squares systems}
  Each pixel $k$ poses the same weighted least-squares problem
  \[
    \hat{x}_k = \arg\min_{x}\,\bigl\lVert W_k^{1/2}\,(G\,x - d_k)\bigr\rVert_2^2 ,
  \]
  with $K$ interferograms (network edges) and $P = D - 1$ unknowns for $D$
  acquisition dates (vertices): the design matrix $G$ is $K \times P$ and
  shared by every pixel, while the phases $d_k$ and the coherence-based
  weights $W_k$ change from pixel to pixel. A scene has $n = 10^{5}$ to
  $10^{8}$ pixels: 3.4~million on the Gal\'apagos scene, about
  $3 \times 10^{8}$ per NISAR frame.

  \centering
  \includegraphics[width=0.95\linewidth]{tall_skinny_matrix.pdf}
  \figcaption{One pixel's system drawn to scale (Fernandina: $K=288$, $P=97$).
    The 191 redundant rows are what make the network well-conditioned (block~5).}

  \raggedright
  Per-pixel cost is $O(KP^2 + P^3/3)$. As the archive lengthens,
  $K \propto D$ and $P^2 \propto D^2$, so the total work grows like $D^3$:
  one to three years of data is about $29\times$ the work.
\end{block}

\end{column}

\begin{column}{\gutter}\end{column}

% ===========================================================================
% COLUMN 2 -- method
% ===========================================================================
\begin{column}{\colw}

% Block 4: the method, framed by the failed direct port. 1.43x is Phase 1
% (report_torch.md); 3,841,835 launches per chunk of 19,403 pixels and 16.5x
% are S1 (report_solver_comparison.md). Both ratios are MintPy-reported step
% times with interpreter start-up excluded, NOT solver-only numbers. Never
% quote the S1 step wall (4.49x): ~47 s of start-up sits on both sides.
% Digest only: no code box and no implementation tips (chunking, rank flags,
% opt-in key); those live in docs/qa_prep.md and on GitHub (block 7's QR codes).
\begin{block}{4 \ Method: reformulate for the GPU, don't just port}
  \keypoint{Porting the CPU solver to a GPU is not enough: the problem has to
    be recast into a shape the GPU is good at.}

  \textbf{A direct port falls short.}
  Replacing the per-pixel \texttt{lstsq} with a batched QR solver on the GPU
  reached only \alert{1.43$\times$} over the CPU path (Fernandina).
  QR sweeps the $P$ columns of every matrix one after another, so a chunk of
  about 19{,}400 pixels issued \alert{3{,}841{,}835 kernel launches}, and a
  GPU is at its worst on many tiny, ordered operations.

  \textbf{Recast around the normal equations,}
  \[
    N_k\,x_k = r_k, \qquad N_k = G^{\!\top} W_k\, G, \quad r_k = G^{\!\top} W_k\, d_k ,
  \]
  where $N_k$ is a small $P \times P$ symmetric positive-definite block.
  In exact arithmetic this is the same least-squares solution, but now every
  $N_k$ of a chunk is built by one batched matrix product (GEMM) and
  factorized by one batched Cholesky call, in single precision with
  PyTorch~[5] and no custom GPU code.
  Missing interferograms enter as zero weights, which is exactly row
  deletion, so every pixel keeps the same shape and a chunk stays one batch.

  \centering
  \includegraphics[width=0.98\linewidth]{fig1_batched_cholesky.pdf}
  \figcaption{The same chunk on both paths: QR keeps $P$ dependent steps
    inside every matrix, while the normal-equation path is a handful of
    batched calls. Kernel launches per chunk measured on Fernandina
    ($K=288$, $P=97$).}

  \raggedright
  \textbf{Same GPU, same scene (Fernandina):} QR 228~s $\rightarrow$ Cholesky
  13.8~s, \alert{16.5$\times$}. Both ratios in this block are step times
  reported by MintPy, with interpreter start-up excluded.
  Gram matrix plus Cholesky is an established strategy for
  parallel and GPU settings, at the documented cost of squaring the condition
  number~[3,\,4]; here only the solution is computed, never a QR factor.
\end{block}

% Block 5: accuracy. Logic for a non-specialist: what kappa means -> why the
% normal equations are the textbook bad choice -> the cost matters only if it
% reaches what InSAR resolves -> measured. Table: report_condition_number.md
% @ d74df50. Error: S2 Galapagos diff JSON (per-pixel RMS vs CPU, median
% 0.065 um, max 16.4 um). Gate: S3 numerical section (5/5 on final
% products). This block absorbed the former "numerical agreement" block; the
% radar-coordinate caveat and CholeskyQR2 are left to docs/qa_prep.md.
% The CPU reference is scipy.linalg.lstsq (default driver gelsd, SVD-based).
\begin{block}{5 \ The accuracy cost: micrometres, against millimetre signals}
  The \emph{condition number} $\kappa$ of a matrix measures how strongly it can
  amplify small errors, such as rounding, on their way into the solution:
  $\kappa \approx 1$ passes them through, a large $\kappa$ magnifies them.

  \centering
  \includegraphics[width=0.85\linewidth]{network_redundancy_conditioning.pdf}
  \figcaption{$\kappa(G)$ is how unevenly $G$ stretches different directions.
    A redundant SBAS network supports every direction and keeps $\kappa$ small.}

  \raggedright
  Forming the normal equations squares it, $\kappa(N_k) = \kappa(G_w)^2$ with
  $G_w = W_k^{1/2} G$. That is why textbooks solve least squares with QR or
  SVD, as the CPU path does, and in single precision Cholesky fails outright once $\kappa(G_w)$ nears
  $4 \times 10^{3}$. The cost is worth paying only if the extra error stays
  below what InSAR can resolve, so we measured both the conditioning of real
  networks and the error of the result.

  \centering
  {\small
  \begin{tabular}{lrrr}
    \toprule
    Scene & $\kappa$, unweighted & $\kappa$, worst weighted pixel & margin \\
    \midrule
    Fernandina        & 11.0 & 323 & $12.7\times$ \\
    Gal\'apagos       & 12.7 & 242 & $16.9\times$ \\
    Kuju              & 28.6 & 81  & $50\times$   \\
    San Francisco Bay & 55.1 & 721 & $5.7\times$  \\
    San Francisco     & 51.1 & (run unweighted) & $80\times$ \\
    \bottomrule
  \end{tabular}\par}
  \figcaption{Float64 SVD on 2{,}100 sampled pixels per scene. Margin is the
    distance from the worst $\kappa$ of the system actually solved to the
    breakdown point, $4.1 \times 10^{3}$.}

  \raggedright
  Against the CPU reference, the per-pixel RMS difference of the time series
  is 0.06~$\mu$m for the median pixel and 16~$\mu$m for the worst one
  (Gal\'apagos, 3.4~million pixels).
  All five scenes pass an agreement gate, RMS difference below $10^{-5}$ of
  the largest CPU value, on the final products (velocity and geocoded time
  series).
  The few pixels that would break down are low-coherence ones that the
  standard temporal-coherence mask discards anyway.

  \centering
  \includegraphics[width=0.8\linewidth]{error_budget.pdf}

  \keypoint{Against the CPU's SVD-based solver, squaring $\kappa$ does leave a
    trace, but at micrometres, buried under InSAR's other error sources.}
\end{block}

\end{column}

\begin{column}{\gutter}\end{column}

% ===========================================================================
% COLUMN 3 -- results and significance
% ===========================================================================
\begin{column}{\colw}

% Block 6: speed. The chart is the main claim. Big number = S2 internal
% (MintPy-timed, start-up excluded: 6169.3 s / 139.0 s); the S2 step wall
% (36.4x) carries ~20-30 s of start-up and is not quoted. Chart = S3 with
% rows labelled by pixels actually inverted, so its Galapagos (1.29 M px after
% masking) is not mistaken for the 3.4 M-pixel S2 run. No size dependence!
% DEM-error (6.2x) and the 18-step run live in docs/qa_prep.md Q4 and the
% caption.
\begin{block}{6 \ Results: speed}
  Every speedup is measured against the CPU reference path on one consumer
  workstation GPU (NVIDIA GeForce RTX~5080, 16~GB).

  \bignumber{6{,}169 s $\rightarrow$ 139 s}{Gal\'apagos scene, 3.4~million
    pixels $\times$ 475 retained interferograms: \textbf{44.4$\times$} as
    timed by MintPy (interpreter start-up excluded), within 7.6~GB of GPU
    memory}

  \centering
  \includegraphics[width=0.97\linewidth]{fig2_speedup.pdf}
  \figcaption{Wall time of the network-inversion step in complete 18-step
    runs on five public scenes (four processors, two radar wavelengths, all on
    Zenodo), with the pixels left after masking. On the log axis the arrow
    length is the speedup. No size dependence is claimed: the ranking does not
    follow pixel count.}
\end{block}

% Block 7: why it matters -- the main push of the closing block. Data growth
% (NISAR, commercial small-SAR constellations, Taiwan's FORMOSAT-9) makes
% consumer-GPU processing matter. FORMOSAT-9 stays regional context with no
% resolution figure. No NISAR extrapolation numbers here (repetition of block
% 3); they live in docs/compute_scale.md and docs/qa_prep.md. Usage details
% (template keys, install, CPU path unchanged) are left to GitHub via the QR
% codes. PR status checked 2026-10-07 -- re-check before printing.
\begin{block}{7 \ Why it matters: SAR data is multiplying}
  \begin{itemize}
    \item NISAR~[1] maps nearly all of Earth's land and ice every 12~days.
    \item Commercial constellations of small SAR satellites add dozens more.
    \item Taiwan's own SAR satellites, FORMOSAT-9~[6], are planned.
  \end{itemize}
  More acquisitions mean longer archives, whose inversion cost grows like
  $D^3$ (block~3), and every new acquisition means inverting the whole network
  again.
  Turning that step from hours on a CPU into minutes on one consumer GPU makes
  routine wide-area monitoring and archive reprocessing practical without an
  HPC cluster.

  \keypoint{As SAR archives multiply, a consumer GPU keeps their time-series
    analysis within reach of any research group.}

  The solver is open source and under review for MintPy (pull request
  \#1490~[7]).

  \begin{columns}[t]
    \begin{column}{0.48\linewidth}
      \centering
      \qrcode[height=0.55\linewidth]{\prurl}\\[0.3ex]
      {\small Code: PR \#1490}
    \end{column}
    \begin{column}{0.48\linewidth}
      \centering
      \qrcode[height=0.55\linewidth]{\posterurl}\\[0.3ex]
      {\small Poster and benchmarks}
    \end{column}
  \end{columns}
\end{block}

% References: same seven entries and numbering as the accepted abstract.
\begin{block}{References}
  \scriptsize
  \begin{enumerate}
    \item NASA-ISRO SAR (NISAR) Mission Science Users' Handbook, NASA JPL, 2018.
    \item Z.~Yunjun, H.~Fattahi, F.~Amelung, ``Small baseline InSAR time series
      analysis: Unwrapping error correction and noise reduction,''
      \emph{Comput.\ Geosci.}, 133, 104331, 2019.
    \item Y.~Yamamoto, Y.~Nakatsukasa, Y.~Yanagisawa, T.~Fukaya, ``Roundoff error
      analysis of the CholeskyQR2 algorithm,'' \emph{ETNA}, 44, 306--326, 2015.
    \item T.~Fukaya, Y.~Nakatsukasa, Y.~Yanagisawa, Y.~Yamamoto, ``CholeskyQR2:
      A simple and communication-avoiding algorithm for computing a tall-skinny
      QR factorization on a large-scale parallel system,'' METR 2014--37,
      The University of Tokyo, 2014.
    \item A.~Paszke et al., ``PyTorch: An imperative style, high-performance
      deep learning library,'' \emph{NeurIPS}, 32, 2019.
    \item Taiwan Space Agency (TASA), FORMOSAT-9 mission,
      \texttt{tasa.org.tw/en-US/missions/detail/FORMOSAT-9}.
    \item S.~Sasaki, ``Add opt-in torch GPU solver for invert\_network,''
      insarlab/MintPy, Pull Request \#1490, 2026 (under review).
  \end{enumerate}
\end{block}

\end{column}
\end{columns}
\end{frame}
\end{document}
